\documentclass{article}
\usepackage{mathnotes}

\title{Joint Probability}
\description{Bivariate random variables, marginalization, multinomial distributions, convolutions, and covariance-correlation analysis.}

\begin{document}

\section{Joint Probability}

\begin{definition}[Joint Probability Distribution]
Two \@{random-variables} $X$ and $Y$ can be paired as a single \textbf{random vector} or \textbf{bivariate random variable.} If it's discrete (i.e. both $X$ and $Y$ are discrete) it has a \textbf{joint probability mass function} and if it's continuous ($X$ and $Y$ are both continuous) it has a \textbf{joint probability density function.}

In the discrete case, $f(x,y) = P(X = x, Y = y).$
\end{definition}

\begin{definition}[Marginalization]
We can \textbf{marginalize} over $Y$ to find solely the distribution of $X,$ we denote this $f_X(x)$ and give it as

\[ f_X(x) = \sum_y f(x, y). \]

We can similarly find $f_Y(y)$ by marginalizing over $X$:

\[ f_Y(y) = \sum_x f(x, y). \]
\end{definition}

\subsection{Multinomial Distribution}

\begin{definition}[Multinomial Distribution]
For a sequence of $n$ independent, identical experiments with each one of the experiments resulting in $r$ outcomes with probabilities $p_1, p_2, \dots, p_r,$ respectively, where $\sum_{i=1}^r p_i = 1,$ we let $X_i$ count the number of the $n$ experiments that result in the $i$th of the $r$ outcomes. Then

\[ P(X_1 = n_1, X_2 = n_2, \dots, X_r = n_r) = \binom{n}{n_1, n_2, \dots, n_r}p_1^{n_1}p_2^{n_2}\cdots p_r^{n_r}, \]

where $\sum_{i=1}^r n_i = n.$
\end{definition}

\begin{definition}[Multinomial Coefficient]
The \@{multinomial-distribution} uses the \textbf{multinomial coefficient}, which is defined as

\[ \binom{n}{n_1,n_2,\dots,n_r} = \frac{n!}{n_1!n_2! \cdots n_r!}. \]
\end{definition}

\subsection{Sums of Independent Random Variables}

\begin{definition}[Convolution]
Sums of independent random variables are called \textbf{convolutions}.

If $X$ and $Y$ are continuous, the probability density function of $X + Y$ is given by

\[ f_{X+Y}(x) = \int_{-\infty}^{\infty} f_X(t) f_Y(x - t) dt. \]

If they are discrete, the probability mass function of $X + Y$ is given by

\[ f_{X+Y}(n) = \sum_k f_X(k)f_Y(n - k). \]
\end{definition}

\subsubsection{Sums of Independent Poisson Random Variables}

\begin{theorem}[Sum of Independent Poisson Random Variables]\label{sum-independent-poisson}
If $X_1, X_2, \dots X_n$ are independent \@[Poisson]{poisson-distribution} random variables with parameters $\lambda_1, \lambda_2, \dots, \lambda_n,$ then $X_1 + X_2 + \cdots + X_n$ is a Poisson random variable with parameter $\lambda_1 + \lambda_2 + \cdots  + \lambda_n.$
\end{theorem}

\subsubsection{Sums of Independent Normal Random Variables}

\begin{theorem}[Sum of Independent Normal Random Variables]\label{sum-independent-normal}
If $X_1, X_2, \dots X_n$ are independent \@[normal]{normal-distribution} random variables with means $\mu_1, \dots, \mu_n$ and \@{variances} $\sigma_1^2, \dots, \sigma_n^2,$ then the random variable $X_1 + X_2 + \cdots + X_n$ is normal with mean $\mu_1 + \mu_2 + \cdots + \mu_n$ and variance $\sigma_1^2 + \sigma_2^2 + \cdots + \sigma_n^2.$
\end{theorem}

\subsection{Covariance and Correlation Coefficient}
\begin{definition}[Covariance]
\textbf{Covariance} is a measure of the joint variability of two \@{random-variables}.

\[ \Cov(X,Y) = E[(X - E[X])(Y - E[Y])]. \]

If large $X$ values go with large $Y$ values and small $X$ goes with small $Y$, covariance will be positive. The covariance will be negative if large $X$ values go with small $Y$ values and vice-versa. An alternative, equivalent form of covariance is:

\[ \Cov(X,Y) = E[XY] - E[X]E[Y]. \]
\end{definition}

\begin{definition}[Correlation Coefficient]
The \textbf{correlation coefficient} of $X$ and $Y$, $\rho_{x,y},$ is given by

\[ \rho_{x,y} = \frac{\Cov{(X, Y)}}{\sigma_X \sigma_Y}. \]
\end{definition}

\end{document}
